\section{Introduction}
\label{sec:introduction}

\subsection{From Search Comparison to Full Portfolio}

G.14~\citep{dellaLibera2026g14} compared eight search-layer algorithms in the
\texttt{redist} ensemble platform --- Short-Burst, SMC, Flip, Forest ReCom,
Merge-Split, Multi-scale MCMC, ShortBurstForest, and VRA-Aware ReCom --- on
three states, four metrics.
That comparison answered a narrow but pressing question: given a fixed initial
plan and a fixed graph, which search algorithm should I run?

The three-layer compositor (\texttt{--structure}, \texttt{--weights-override},
\texttt{--search}) makes this narrow framing incomplete.
A practitioner who chooses \textsc{Multi-scale} as the search algorithm but
uses the default METIS bisection as the structure algorithm may obtain a
plan that is mixing efficiently but starting from a suboptimal compactness
baseline.
Conversely, a practitioner who runs Simulated Annealing structure (best
heuristic \ec{} in the structure layer) with Flip search (local sensitivity
only) is using an expensive structure algorithm to build a starting point
that the search algorithm will never meaningfully improve.

This paper provides the first systematic comparison of all 15 implemented
algorithms across both structure and search layers.
Our goal is practical: give a practitioner a complete reference for the full
algorithm portfolio.

\subsection{The 15 Algorithms}

The \texttt{redist} platform implements six structure algorithms and twelve
search algorithms, totalling 15 distinct algorithm identities (three search
algorithms are new since G.14).

\paragraph{Structure layer.}
Structure algorithms execute one full bisection of the tract graph to produce
an initial $k$-district plan.
The five non-ILP algorithms are applied recursively; ILP is applied once per
bisection node.

\begin{enumerate}
  \item \textbf{Standard METIS} (\metis{}, baseline): METIS multilevel
        $k$-way partitioning~\citep{karypis1998}.
        Fast ($O(N \log N)$), well-tested, default structure algorithm.

  \item \textbf{Simulated Annealing} (\sa{}, B.19): METIS initialisation
        followed by SA improvement steps on the bisection
        graph~\citep{dellaLibera2026sa}.
        Best \ec{} among heuristic structure methods; 12$\times$ slower than
        METIS.

  \item \textbf{CVD Graph-Distance} (\cvdg{}, B.22): Centroidal Voronoi
        Diagram bisection using graph hop-distance as the metric, requiring
        no geographic data beyond the
        adjacency graph~\citep{dellaLibera2026cvd}.
        Best \pp{} among structure methods when centroids are unavailable.

  \item \textbf{CVD Geographic} (\cvdgeo{}, B.22b): CVD bisection using
        Euclidean distance from tract centroids~\citep{dellaLibera2026cvdgeo}.
        Best \pp{} among all structure methods; requires centroid data.

  \item \textbf{BFS Growth} (\bfsgrow{}, B.23): Breadth-first search growth
        from $k$ seeds, fastest structure
        method~\citep{dellaLibera2026bfsgrowth}.
        Zero hyperparameters, $O(N \log N)$; \pp{} competitive with \cvdg{}.

  \item \textbf{ILP Exact} (\ilp{}, B.24): Integer linear program solving
        the minimum-\ec{} bisection exactly for small
        instances~\citep{dellaLibera2026ilp}.
        Provides certified \ec{}-optimality; applicable only when $N \leq 500$
        (NH $k = 2$, each half $\leq 300$ tracts).
\end{enumerate}

\paragraph{Search layer.}
Eight algorithms were compared in G.14; this paper adds three new algorithms
(marked with $\star$):

\begin{enumerate}
  \item \textbf{ConvergenceSweep} (\cs{}, B.16): Multi-seed convergence
        sweep; federal statute default~\citep{dellaLibera2026convSweep}.
  \item \textbf{PercentileSweep} (\ps{}, H.0): Statutory choice of legal
        posture via percentile
        selection~\citep{dellaLibera2026percentileSweep}.
  \item \textbf{BisectionEnsemble} (\be{}, H.1): Local ReCom at bisection
        nodes~\citep{dellaLibera2026bisectionEns}.
  \item \textbf{Short-Burst} (\sburst{}, G.6)~\citep{dellaLibera2026shortburst}.
  \item \textbf{Flip} (\flip{}, G.8)~\citep{dellaLibera2026flip}.
  \item \textbf{Forest ReCom} (\fr{}, G.9)~\citep{dellaLibera2026forestRecom}.
  \item \textbf{Merge-Split} (\msp{}, G.10)~\citep{dellaLibera2026mergeSplit}.
  \item \textbf{Multi-scale} (\ms{}, G.11)~\citep{dellaLibera2026multiscale}.
  \item \textbf{ShortBurstForest} (\sburstf{}, G.12)~\citep{dellaLibera2026shortburstChains}.
  \item \textbf{VRA-Aware ReCom} (\vra{}, G.13)~\citep{dellaLibera2026vraAware}.
  \item $\star$ \textbf{Parallel Tempering} (\pt{}, B.20): Replica exchange
        across temperature ladder; best \ec{} for large
        $k$~\citep{dellaLibera2026pt}.
  \item $\star$ \textbf{SMC-Percentile} (\smcp{}, G.7/new): Calibrated SMC
        selecting the median plan by \pp{}~\citep{dellaLibera2026smcPercentile}.
  \item $\star$ \textbf{Adaptive Multi-scale} (\ams{}, B.21): Self-tuning
        coarse fraction $\alpha$~\citep{dellaLibera2026adaptiveMS}.
\end{enumerate}

\subsection{Five Test States}

G.14 used NC ($k = 14$), WI ($k = 8$), and TX ($k = 38$).
G.15 adds:

\begin{itemize}
  \item \textbf{Florida (FL)}, $k = 28$.
        A large-$k$ state between TX and NC in district count, testing
        whether the TX findings for Parallel Tempering and Adaptive Multi-scale
        generalise to an intermediate $k$.
        FL has $N \approx 4{,}100$ census tracts.

  \item \textbf{New Hampshire (NH)}, $k = 2$.
        The only state for which ILP is tractable: two congressional districts
        with $N \approx 260$ tracts.
        NH provides the \textit{certified-optimal} baseline for all structure
        algorithms: ILP computes the provably minimum-\ec{} bisection.
\end{itemize}

\subsection{Main Findings}

Five practical conclusions emerge from the full portfolio comparison.

\begin{enumerate}

  \item \textbf{ILP certifies what heuristics approximate.}
        For NH ($k = 2$), \ilp{} achieves \ec{} = 47 (certified optimal).
        SA achieves 51 (nearest heuristic), CVD-Geo 53, BFS 56, METIS 58.
        The optimality gap of the best heuristic is 8.5\%; the METIS gap
        is 23.4\%.

  \item \textbf{Structure layer determines the PP ceiling.}
        \cvdgeo{} achieves the highest \pp{} among all 15 algorithms across
        all states, outperforming even \smcp{} (the best search method for
        \pp{}).
        PP is largely determined at bisection time; search algorithms raise
        PP modestly from the structure baseline.

  \item \textbf{Parallel Tempering dominates for large $k$.}
        For TX ($k = 38$) and FL ($k = 28$), \pt{} achieves lower final
        \ec{} than \sburst{} by escaping basin boundaries that
        burst strategies cannot cross without
        the temperature exchange mechanism.
        For small $k$ (WI, NC), \sburst{} remains competitive.

  \item \textbf{Adaptive Multi-scale is superior to fixed Multi-scale for
        TX.}
        \ams{} auto-tunes the coarse fraction $\alpha$ in real time,
        allocating more coarse steps during the warm-up phase when
        the chain is far from the mode.
        For TX, \ams{} reduces \acf(100) to 0.31 vs.\ 0.37 for fixed
        \ms{}; for NC, they are indistinguishable.

  \item \textbf{The full decision framework requires six questions, not four.}
        The G.14 four-question tree selects only the search algorithm.
        Two additional questions --- (v) Is exact structure optimality
        required? (vi) Is geographic compactness (\pp{}) the primary
        structure criterion? --- are needed to select the structure
        algorithm.
        The expanded framework covers all 15 algorithms.

\end{enumerate}

\subsection{Relation to G.14 and Prior Work}

G.14~\citep{dellaLibera2026g14} is the immediate predecessor.
Results from G.14 for the original eight search algorithms on NC/WI/TX are
carried forward without re-running; new states (FL, NH) and new algorithms
(\pt{}, \smcp{}, \ams{}) are added with dagger notation indicating new
single-run estimates.
Structure-layer results are entirely new; G.14 fixed structure at
\texttt{ratio-optimal} for all comparisons.

Cannon et al.~\citep{cannon2022} and Janson~\citep{janson2022} provide
theoretical foundations for short-burst and chain mixing, respectively, but
do not address structure algorithms.
This paper provides the first joint structure+search comparison in the
redistricting literature.

\subsection{Paper Structure}

Section~\ref{sec:setup} extends the G.14 experimental setup to five states
and 15 algorithms.
Section~\ref{sec:structure} compares the six structure algorithms.
Section~\ref{sec:search} extends the G.14 search comparison with the three
new algorithms.
Section~\ref{sec:distributional} examines distributional properties including
\pt{} and \ams{}.
Section~\ref{sec:decision} presents the six-question decision framework and
combined structure+search recommendation table.
Section~\ref{sec:budget} gives compute budget guidance for a 50-state
production sweep.
Section~\ref{sec:conclusion} concludes.
